\chapter{Discussion and conclusion}\label{conclusions}
This study investigated the roles that \acrshort{gpt} plays in conversations with software developers, using a large, multi-source dataset of developer-\acrshort{gpt} interactions. Using a combination of n-gram analysis, \acrfull{spm}, and topic modeling across seven distinct data sources, which include GitHub commits, issues, discussions, pull requests, code files, repositories, and HackerNews, we studied the lexical and thematic structures that characterize developer expectations and chatbot functionalities.

Key findings show contextual dependence of \acrshort{gpt} roles, derived by the origin of the conversation. For example, GitHub commits mainly reflect task automation and file-based scripting workflows, which supports the view that developers treat \acrshort{gpt} as a procedural collaborator, trusted with command-line logic, automated code generation, and low-level implementation tasks. These findings are in-line with the work of Sagdic et al.~\cite{sagdic2024-devgpt-taxonomy} and Mohamed et al.~\cite{Mohamed2024-topics}, who also observed dominant use cases in code generation and automation.

Contrarily, conversations sourced from GitHub issues and discussions show a different interaction style. Topics in these sources express debugging logic, system configuration, and architectural reasoning. Discussion threads reflected higher-level planning and storytelling, with themes like \textbf{Product design feasibility} and \textbf{Silicon Valley series} showing that \acrshort{gpt} is sometimes treated as a reasoning or storytelling partner, and not only as a programming assistant.

Conversations from GitHub pull requests and code files illustrated the importance of semantic and syntactic correctness. Developers often use \acrshort{gpt} for tasks involving function validation, API usage, and type system clarification. These contexts demand higher trust in the chatbot's technical accuracy, yet the literature cautions against full automation due to persisting issues with correctness and code safety~\cite{siddiq2023generate, gpt-code-security-quality, clark2024-quantitative-quality-analysis}.

Finally, conversations sourced from GitHub repositories and HackerNews revealed \acrshort{gpt}'s capacity to support documentation, vision statements, and informal discourse. Sequences like \textit{``fill height viewport leave position foreground left edge right''} suggest that \acrshort{gpt} is used to assist in UI layout composition and document structuring. Meanwhile, HackerNews conversations uncovered more dialogic and playful use cases, showing \acrshort{gpt} serving as a narrative generator, ideation tool, and storytelling agent. This supports previous research emphasizing the model's capabilities in non-code domains~\cite{fi15060192}.

From a methodological perspective, the integration of \acrshort{spm} and topic modeling yielded some interpretable patterns. The findings demonstrate how varied \acrshort{gpt}'s role is in software development, ranging across domains from automation and testing to planning and narrative creation. These findings contribute to the research question posed in this study by showing that \acrshort{gpt}'s perceived role spans over code assistant, design partner, documentation helper, and conversational companion.

\section{Future work}
Future research could extend this analysis by incorporating \acrshort{gpt} responses, allowing for a more holistic understanding of co-construction and feedback loops in developer-chatbot dialogues. Additionally, fine-grained intent classification and temporal analysis could reveal evolving usage trends and adaptation behaviors. Lastly, there is potential to design tailored \acrlong{llm}s or prompt-engineering guidelines optimized for the distinct roles developers expect \acrshort{gpt} or other chatbots to fulfill.